\section{Results}\label{sec:results}

\subsection{Information structure beyond standard factors}
\label{sec:results-factor-value}

The information field measured from news organizes residual return dependence
in the later evaluation period.
With all parameters fitted on 2018--2022 returns, the barycentric field
$W^\flat$ raises the
2023--2026 mean conditional quasi-log score by
\ppnum[4]{spatial-value-spherical-w2-versus-none-held-out-score-gain} per
firm-day relative to the factor-only model, with 95\% paired interval
$[\ppnum[4]{spatial-value-spherical-w2-versus-none-held-out-score-gain-ci-lower},
\ppnum[4]{spatial-value-spherical-w2-versus-none-held-out-score-gain-ci-upper}]$.
The positive gain means that information-peer relations improve the joint
distribution assigned to subsequent returns after conditioning on standard
factors.
The linear-cost W1 field also improves the score, with its
interval above zero.

\Cref{tab:spatial-value} compares the fields with the matched factor-only
model, holding FF5-plus-momentum controls, firm coefficients, dates, and the
innovation specification fixed.
The field adds a single spatial coefficient.
On the post-corpus estimation panel, W1 gains
\ppnum[1]{spatial-value-spherical-w1-versus-none-loglik-gain} QLL units and W2
gains \ppnum[1]{spatial-value-spherical-w2-versus-none-loglik-gain}, and the
restricted-null bootstrap rejects zero spatial feedback for both fields, at
$p=\ppnum[4]{spatial-value-spherical-w1-versus-none-qlr-pvalue}$ for W1 and
$p=\ppnum[4]{spatial-value-spherical-w2-versus-none-qlr-pvalue}$ for W2.
The information criteria also favor each spatial model
(\ppAppendixRef{sec:app-information-criteria}).

\begin{table}[H]
  \centering
  \caption{Primary spatial-value comparison under spherical innovations.
    Every row uses FF5 plus momentum and the same firm-specific coefficients.
    QLL is maximized on the post-corpus panel. Restricted-null QLR $p$-values
    use 2,000 independently resampled firm paths under the working $\rho=0$
    null; zero exceedances are reported at the add-one floor.
    Out-of-period conditional score gains are relative to the factor-only model,
    with 95\% paired stationary-bootstrap intervals; higher is better.
  Evaluation-date factors are observed when scoring.}
  \label{tab:spatial-value}
  \begin{tabular}{lrrrr}
    \toprule
    Model & $\hat\rho$ & QLL & QLR $p$ & Held-out gain \\
    \midrule
    Factor only &
    \ppnum[3]{spatial-value-spherical-none-rho} &
    \ppnum[1]{spatial-value-spherical-none-loglik} &
    --- & --- \\
    Target-anchored W1 &
    \ppnum[3]{spatial-value-spherical-w1-rho} &
    \ppnum[1]{spatial-value-spherical-w1-loglik} &
    \ppnum[4]{spatial-value-spherical-w1-versus-none-qlr-pvalue} &
    $\ppnum[4]{spatial-value-spherical-w1-versus-none-held-out-score-gain}$
    $[\ppnum[4]{spatial-value-spherical-w1-versus-none-held-out-score-gain-ci-lower},
    \ppnum[4]{spatial-value-spherical-w1-versus-none-held-out-score-gain-ci-upper}]$
    \\
    Target-anchored W2 &
    \ppnum[3]{spatial-value-spherical-w2-rho} &
    \ppnum[1]{spatial-value-spherical-w2-loglik} &
    \ppnum[4]{spatial-value-spherical-w2-versus-none-qlr-pvalue} &
    $\ppnum[4]{spatial-value-spherical-w2-versus-none-held-out-score-gain}$
    $[\ppnum[4]{spatial-value-spherical-w2-versus-none-held-out-score-gain-ci-lower},
    \ppnum[4]{spatial-value-spherical-w2-versus-none-held-out-score-gain-ci-upper}]$
    \\
    \bottomrule
  \end{tabular}
\end{table}

For the primary field $W^\flat$, the factor-conditioned estimate is
$\hat\rho=\ppnum[3]{transport-w2-rho}$, with 95\% interval
$[\ppnum[3]{transport-w2-rho-ci-lower},
\ppnum[3]{transport-w2-rho-ci-upper}]$, and the corresponding
adjustment index is
$\hat\lambda=\ppnum[2]{transport-w2-lambda}$, with interval
$[\ppnum[2]{transport-w2-lambda-ci-lower},
\ppnum[2]{transport-w2-lambda-ci-upper}]$.
Factor conditioning matters for magnitude: the matched raw-return,
common-intercept estimate is $\ppnum[3]{transport-w2-unadjusted-rho}$, so the
factors absorb part, but only part, of the dependence that a raw-return model
assigns to the field.
The remaining dependence is stable across common-factor controls: the
estimate ranges from \ppnum[3]{transport-factor-ff3-w2-rho} under FF3 to
\ppnum[3]{transport-factor-ff5-mom-w2-rho} under FF5 plus momentum, and every
bootstrap interval excludes zero
(\ppAppendixRef{sec:app-information-criteria}).
That stability parallels the co-mention evidence of
\citet{ge_news-implied_2023}, whose firm coefficients and sample differ from
ours.

The improvement concerns how unexpected returns move together.
Under the unrestricted firm coefficients of
\Cref{prop:factor-beta-invariance}, reduced-form factor betas and fitted means
equal their firm-by-firm least-squares values; the estimated models reproduce
that identity to within $10^{-10}$.
The held-out gain therefore reflects residual covariance.
Allowing every firm its own innovation variance raises the estimates to
\ppnum[3]{spatial-value-diagonal-w1-rho} for W1 and
\ppnum[3]{spatial-value-diagonal-w2-rho} for W2, rejects $\rho=0$ for both
fields, and keeps both held-out intervals above zero
(\Cref{tab:spatial-value-diagonal}).

\subsection{Anatomy of the barycentric field}
\label{sec:results-anatomy}

\Cref{fig:peer-basket} opens one row of the estimated field.
AMD's basket combines firms from complementary chip-design, equipment-supply,
and memory segments, led by NVDA, INTC, MU, AMAT, and ASML.
No single peer dominates: the five largest weights account for less than half
of the row, and the remaining mass is spread across many firms.
The row is a synthetic information-peer benchmark for AMD, constructed without
reference to industry codes or returns.

\begin{figure}[htbp]
  \centering
  \ppfigure{1}{peer_basket}
  \caption{Anatomy of one barycentric peer basket. Bars show AMD's five largest
    weights $W^\flat_{ij}$ in the estimated field and the combined weight of
    all remaining peers; the row sums to one and excludes AMD. The right-hand
    column gives
  each peer's industry; peers in AMD's industry are set in black.}
  \label{fig:peer-basket}
  \figalttext{Horizontal bar chart for target AMD. Five named peer firms carry
    the largest reconstruction weights, each well below one half, and a final
  bar collects the weight spread over all remaining firms.}
\end{figure}

Reconstruction relevance differs from pairwise proximity.
\Cref{fig:peer-definitions} places one target, HON, under three peer
definitions built from the same 2018--2022 information.
HON is selected mechanically as the target with the largest effective source
count, using only the frozen information field and no return outcomes; its
barycentric row is therefore the most diffuse in the sample, which makes the
contrast sharp. With the median bandwidth, the pairwise RBF row is nearly
uniform over all candidates, tilting only slightly toward the firms nearest
to HON in Wasserstein distance, and the co-mention row concentrates on a few
firms that appear alongside HON in news articles. The barycentric row selects
the combination that jointly reconstructs HON's footprint, and it assigns
UAL, only the 51st-closest firm by
pairwise distance, the second-largest weight
(\Cref{tab:p5-peer-interpretation}).
A firm far from the target on average can still contribute strongly to
reproducing one part of its information distribution, and the full field is
directed and heterogeneous across targets
(\ppAppendixRef{sec:app-field-anatomy}).

\begin{figure}[htbp]
  \centering
  \ppfigure{1}{peer_definitions}
  \caption{One target under three peer definitions. Panel-specific scales are
    used; compare composition, not bar length across panels. Each
    panel shows HON's
    five largest weights and the combined remaining weight under the
    barycentric field $W^\flat$ (left), the median-bandwidth pairwise
    RBF field built
    from the same Wasserstein distances (center), and the persistent news
    co-mention field (right). Shading marks
  a peer that appears in more than one panel.}
  \label{fig:peer-definitions}
  \figalttext{Three horizontal bar charts for target HON. The barycentric
    panel spreads weight over a mixed set of peers, including a firm that is
    far from HON in pairwise distance; the RBF panel is nearly flat; the
  co-mention panel concentrates on a few firms reported with HON in news.}
\end{figure}

\subsection{Joint reconstruction and simpler peer definitions}
\label{sec:results-horse-race}

Joint reconstruction provides the largest held-out gain among the tested
peer-field specifications.
\Cref{tab:p5-economic-field-comparison} holds firms, dates, factor controls,
and the return likelihood fixed while changing the peer definition.
Among peer-field specifications, the two barycentric fields record the largest
held-out gains over the factor-only model,
\ppnum[4]{economic-w1-held-out-gain} for W1 and
\ppnum[4]{economic-w2-held-out-gain} for W2.

% AUTO-GENERATED -- do not edit by hand.
\begin{table}[H]
\centering
\scriptsize
\setlength{\tabcolsep}{3.0pt}
\begin{tabular}{lrrrrr}
\toprule
Dependence specification & $\hat\rho$ & QLL & Held-out score & Gain & 95\% CI \\
\midrule
Factor-only & 0.000 & 205748.3 & 2.468881 & 0.000000 & -- \\
Direct Ledoit--Wolf residual covariance & -- & -- & 2.558644 & 0.089763 & [-0.0028, 0.0903] \\
W1 barycentric & 0.684 & 209488.3 & 2.526065 & 0.057184 & [0.0497, 0.0656] \\
W2 barycentric & 0.687 & 209476.0 & 2.525957 & 0.057076 & [0.0496, 0.0653] \\
Median-bandwidth W2 RBF & 0.467 & 206227.6 & 2.479192 & 0.010311 & [0.0082, 0.0129] \\
Calibration-tuned W2 RBF & 0.572 & 206604.0 & 2.484870 & 0.015989 & [0.0135, 0.0191] \\
Centroid reconstruction & 0.470 & 209343.1 & 2.522691 & 0.053810 & [0.0460, 0.0625] \\
Equal active support & 0.619 & 207216.3 & 2.494105 & 0.025224 & [0.0218, 0.0294] \\
Persistent co-mentions & 0.466 & 209240.3 & 2.519320 & 0.050439 & [0.0428, 0.0586] \\
Industry/sector fallback & 0.229 & 207873.3 & 2.501014 & 0.032133 & [0.0259, 0.0390] \\
\bottomrule
\end{tabular}
\caption{Factor-conditioned peer-field and covariance benchmarks. All rows use the primary FF5-plus-momentum controls, the same firms and dates; fields use firm information dated through 2022 and are held fixed throughout the 2023--2026 evaluation, and covariance inputs are calibrated on 2018--2022 returns. The held-out score is the mean Gaussian quasi-log score per firm-day on 2023--2026; gain is relative to the factor-only fit. The direct Ledoit--Wolf row estimates a conventional factor-residual covariance from the calibration sample and therefore has no spatial coefficient or QLL. Intervals are paired stationary-block bootstrap intervals that refit every comparator on each calibration draw, matching the resampling architecture of \Cref{tab:spatial-value}.}
\label{tab:p5-economic-field-comparison}
\end{table}


The conventional residual-covariance benchmark has the higher point estimate
on the same held-out score: its gain is
\ppnum[4]{economic-ledoit-wolf-held-out-gain}, compared with
\ppnum[4]{economic-w2-held-out-gain} for $W^\flat$.
Its own gain interval spans zero once calibration uncertainty enters, while
the peer-field gain intervals remain positive.
Across the common draws the paired Ledoit--Wolf-minus-barycentric
difference is positive
in a share \ppnum[3]{economic-ledoit-wolf-minus-w2-positive-share} of
replications, with 95\% interval
$[\ppnum[4]{economic-ledoit-wolf-minus-w2-ci-lower},
\ppnum[4]{economic-ledoit-wolf-minus-w2-ci-upper}]$.
That interval includes zero, so the comparison does not establish which model
has the higher expected score.
Direct shrinkage estimates covariance itself; the barycentric construction
supplies an explicit peer structure whose financial content is evaluated by
the same score.

Pairwise proximity captures only a small part of that gain.
The median-bandwidth RBF field uses the same empirical laws and the same
pairwise Wasserstein distances as $W^\flat$ but converts each
distance separately
into a weight; its gain is \ppnum[4]{economic-rbf-median-held-out-gain}.
Choosing the RBF bandwidth on the 2018--2022 validation split raises the gain
only to \ppnum[4]{economic-rbf-tuned-held-out-gain}, and neither RBF interval
overlaps the barycentric intervals.
The large increase from proximity weighting to joint reconstruction locates
most of the field's return content in the joint representation of the target.

The fitted weights matter beyond the peer set they select.
Weighting the active peers of $W^\flat$ equally yields a gain of
\ppnum[4]{economic-equal-support-held-out-gain}, less than half the
barycentric gain, with a disjoint interval.
A centroid reconstruction, which collapses each firm's articles to their mean
embedding before solving the same leave-one-out simplex problem, attains
\ppnum[4]{economic-centroid-held-out-gain}.
Because the two models use the same evaluation dates, the paired
barycentric-minus-
centroid score difference is \ppnum[4]{economic-w2-minus-centroid-estimate}
with 95\% interval
$[\ppnum[4]{economic-w2-minus-centroid-ci-lower},
\ppnum[4]{economic-w2-minus-centroid-ci-upper}]$; $W^\flat$ is higher in a share
\ppnum[3]{economic-w2-minus-centroid-positive-share} of the draws.
The centroid therefore retains most of the field's score gain, while the full
distribution contributes a small but measurable increment beyond
target-specific reconstruction.
Persistent news co-mentions attain \ppnum[4]{economic-co-mentions-held-out-gain}
and the industry/sector fallback
\ppnum[4]{economic-industry-sector-held-out-gain}.
The barycentric point estimates exceed both, and the relation between the
language-model field and explicit news links is taken up directly in
\Cref{sec:field-overlap}.
The spatial coefficient itself does not rank the fields: median-bandwidth RBF
and co-mentions have nearly equal $\hat\rho$ but held-out gains that differ by
a factor of five.
In-sample quasi-log-likelihood in the raw-return, common-intercept
specification orders the barycentric, co-mention, equal-support, and RBF
fields in the same way (\Cref{tab:qmle-pooled}).

\subsection{Language-model peers and explicit news links}
\label{sec:field-overlap}

Information reconstruction and shared news coverage yield different peer
maps.
The barycentric and co-mention fields choose partly overlapping peer sets,
with different weights and related peer-return signals.
\Cref{tab:field-overlap} separates peer selection, cardinal weighting, and the
induced peer-return series, and evaluates overlap against matched-random
baskets because row sparsity sets a floor for peer overlap and common return
variation sets one for signal correlation.

\begin{table}[H]
  \centering
  \caption{Agreement between the frozen barycentric interaction field and the
    persistent co-mention field. The null draws
    \ppvalue{overlap-support-draws}
    fixed-seed baskets carrying $W^\flat$'s own row sparsities; $p$ is the
    share of null draws reaching the observed value. Rows with no mechanical
  floor carry no null.}
  \label{tab:field-overlap}
  \begin{tabular}{llrrr}
    \toprule
    Stage & Statistic & Observed & Null & $p$ \\
    \midrule
    Peer sets & Directed share of $W^\flat$ peers &
    \ppnum[3]{overlap-directed-share} &
    \ppnum[3]{overlap-support-null-share} &
    \ppnum[3]{overlap-support-null-pvalue} \\
    & Rank-matched share &
    \ppnum[3]{overlap-rank-matched-share} & --- & --- \\
    & Jaccard overlap &
    \ppnum[3]{overlap-jaccard-mean} & --- & --- \\
    \midrule
    Weights & Off-diagonal correlation &
    \ppnum[3]{overlap-weight-offdiagonal-pearson} & --- & --- \\
    & Cosine on the support union &
    \ppnum[3]{overlap-weight-union-cosine} & --- & --- \\
    & Rank correlation on common edges &
    \ppnum[3]{overlap-weight-common-spearman} & --- & --- \\
    \midrule
    Induced field & Correlation of $W r_t$ &
    \ppnum[3]{overlap-induced-pearson} &
    \ppnum[3]{overlap-induced-null-pearson} &
    \ppnum[3]{overlap-induced-null-pvalue} \\
    & Coefficient of determination &
    \ppnum[3]{overlap-induced-r-squared} & --- & --- \\
    \bottomrule
  \end{tabular}
\end{table}

The barycentric rows are the denser of the two, carrying
\ppnum[1]{overlap-w-flat-support-size} active peers on average against
\ppnum[1]{overlap-co-mentions-support-size} for the co-mention graph, and
\ppnum[3]{overlap-directed-share} of the barycentric peers are also
co-mention peers.
A random basket with the same row sparsities reaches
\ppnum[3]{overlap-support-null-share}, an excess of
\ppnum[1]{overlap-support-excess-pct} percentage points
($p=\ppnum[3]{overlap-support-null-pvalue}$): the two maps overlap more than
chance but remain distinct.
Their cardinal weights agree only partially, with off-diagonal correlation
\ppnum[3]{overlap-weight-offdiagonal-pearson} and a lower rank correlation of
\ppnum[3]{overlap-weight-common-spearman} over the
\ppnum[1]{overlap-weight-common-count} edges per row they share.
The induced peer-return series correlate more strongly, at
\ppnum[3]{overlap-induced-pearson}, but any two row-stochastic averages of a
co-moving cross-section are correlated: a random basket at $W^\flat$'s
sparsity already reaches \ppnum[3]{overlap-induced-null-pearson}.
The observed correlation exceeds that floor
($p=\ppnum[3]{overlap-induced-null-pvalue}$), and per-firm correlations range
from \ppnum[3]{overlap-induced-asset-min} to
\ppnum[3]{overlap-induced-asset-max}.
Common return variation therefore makes the fields look more alike as return
signals than as peer maps.

\subsection{Two channels in one adjustment problem}
\label{sec:two-field-results}

Combining the information field with explicit news links improves held-out
joint density beyond either field alone.
The joint model of \Cref{thm:two-field-closure} assigns a separate coefficient
to each peer definition.
\Cref{tab:two-field-pooled} reports the post-corpus factor-conditioned
estimates, with each channel's index given by \eqref{eq:channel-inversion}.

\begin{table}[H]
  \centering
  \caption{Pooled 2023--2026 joint barycentric and news-link estimates under
    the primary FF5-plus-momentum specification. The 95\% intervals are
    descriptive quantiles from 2,000 joint-date stationary-bootstrap refits.
    The working-null QLR comparisons use restricted factor-conditioned
  stationary-bootstrap calibrations.}
  \label{tab:two-field-pooled}
  \begin{tabular}{lrrrr}
    \toprule
    Channel & $\hat\rho$ & 95\% interval & $\hat\lambda$ & 95\% interval \\
    \midrule
    Barycentric ($W^\flat$) &
    \ppnum[3]{post-pooled-2023-2026-factor-two-field-rho-b} &
    $[\ppnum[3]{post-pooled-2023-2026-factor-two-field-rho-b-ci-lower},
    \ppnum[3]{post-pooled-2023-2026-factor-two-field-rho-b-ci-upper}]$ &
    \ppnum[2]{post-pooled-2023-2026-factor-two-field-lambda-b} &
    $[\ppnum[2]{post-pooled-2023-2026-factor-two-field-lambda-b-ci-lower},
    \ppnum[2]{post-pooled-2023-2026-factor-two-field-lambda-b-ci-upper}]$ \\
    News-link ($W^{\mathrm{news}}$) &
    \ppnum[3]{post-pooled-2023-2026-factor-two-field-rho-n} &
    $[\ppnum[3]{post-pooled-2023-2026-factor-two-field-rho-n-ci-lower},
    \ppnum[3]{post-pooled-2023-2026-factor-two-field-rho-n-ci-upper}]$ &
    \ppnum[2]{post-pooled-2023-2026-factor-two-field-lambda-n} &
    $[\ppnum[2]{post-pooled-2023-2026-factor-two-field-lambda-n-ci-lower},
    \ppnum[2]{post-pooled-2023-2026-factor-two-field-lambda-n-ci-upper}]$ \\
    Total &
    \ppnum[3]{post-pooled-2023-2026-factor-two-field-rho-total} &
    $[\ppnum[3]{post-pooled-2023-2026-factor-two-field-rho-total-ci-lower},
    \ppnum[3]{post-pooled-2023-2026-factor-two-field-rho-total-ci-upper}]$
    & & \\
    \midrule
    \bottomrule
  \end{tabular}
\end{table}

With parameters fitted on the calibration period, the joint model improves
the 2023--2026 score relative to the barycentric-only model by
\ppnum[4]{post-pooled-2023-2026-factor-two-field-held-out-joint-minus-barycentric-estimate}
with 95\% interval
$[\ppnum[4]{post-pooled-2023-2026-factor-two-field-held-out-joint-minus-barycentric-ci-lower},
\ppnum[4]{post-pooled-2023-2026-factor-two-field-held-out-joint-minus-barycentric-ci-upper}]$,
and relative to the news boundary by
\ppnum[4]{post-pooled-2023-2026-factor-two-field-held-out-joint-minus-news-estimate}
with interval
$[\ppnum[4]{post-pooled-2023-2026-factor-two-field-held-out-joint-minus-news-ci-lower},
\ppnum[4]{post-pooled-2023-2026-factor-two-field-held-out-joint-minus-news-ci-upper}]$.
Both intervals lie above zero, so each field adds out-of-period information
conditional on the other.
The restricted-null tests also reject each exclusion: adding the news-link
field to the barycentric boundary gives
$\mathrm{QLR}^{\mathrm{news}}=\ppnum[1]{post-pooled-2023-2026-factor-two-field-qlr-news-given-barycentric-statistic}$
($p=\ppnum[4]{post-pooled-2023-2026-factor-two-field-qlr-news-given-barycentric-pvalue}$),
and adding the barycentric field to the news boundary gives
$\mathrm{QLR}^\flat=\ppnum[1]{post-pooled-2023-2026-factor-two-field-qlr-barycentric-given-news-statistic}$
($p=\ppnum[4]{post-pooled-2023-2026-factor-two-field-qlr-barycentric-given-news-pvalue}$).
Across \ppvalue{post-pooled-2023-2026-factor-two-field-trading-days} common
dates, the joint model improves on the barycentric-only and news-only
boundaries by
\ppnum[1]{post-pooled-2023-2026-factor-two-field-loglik-gain-w-flat} and
\ppnum[1]{post-pooled-2023-2026-factor-two-field-loglik-gain-co-mentions}
quasi-log-likelihood units, respectively.
The barycentric channel receives the larger share: its index is
$\hat\lambda^\flat=\ppnum[2]{post-pooled-2023-2026-factor-two-field-lambda-b}$,
with
interval
$[\ppnum[2]{post-pooled-2023-2026-factor-two-field-lambda-b-ci-lower},
\ppnum[2]{post-pooled-2023-2026-factor-two-field-lambda-b-ci-upper}]$, and the
news-link index is
$\hat\lambda^{\mathrm{news}}=\ppnum[2]{post-pooled-2023-2026-factor-two-field-lambda-n}$,
with
interval
$[\ppnum[2]{post-pooled-2023-2026-factor-two-field-lambda-n-ci-lower},
\ppnum[2]{post-pooled-2023-2026-factor-two-field-lambda-n-ci-upper}]$.
Under the adjustment interpretation of \eqref{eq:two-field-objective}, the
larger barycentric index corresponds to a stronger penalty on misalignment
with the barycentric peer average than with co-mentioned firms.

The two channel coefficients share the total dependence assigned to peers.
The joint total $\hat\rho^\flat+\hat\rho^{\mathrm{news}}=
\ppnum[3]{post-pooled-2023-2026-factor-two-field-rho-total}$ has interval
$[\ppnum[3]{post-pooled-2023-2026-factor-two-field-rho-total-ci-lower},
\ppnum[3]{post-pooled-2023-2026-factor-two-field-rho-total-ci-upper}]$.
The channel coefficients have bootstrap correlation
\ppnum[3]{post-pooled-2023-2026-factor-two-field-channel-correlation},
so total peer
adjustment is estimated more sharply than its division between channels
under the factor-conditioned specification.
Language-model reconstruction and explicit co-coverage are thus related but
economically distinct definitions of a firm's information peers, and the
barycentric definition receives the larger share of adjustment.

The raw-return, common-intercept joint estimates and their uncertainty region
remain in \ppAppendixRef{sec:app-unadjusted} as a corroborating specification.

\subsection{Transport-order equivalence}
\label{sec:results-transport-order}

The financial evidence survives replacing quadratic transport costs with
linear costs.
\Cref{tab:transport-order} compares target-anchored W1 and W2 fields on the
same \ppvalue{transport-n-assets} firms and
\ppvalue{transport-evaluation-trading-days} dates; the two constructions share
the corpus, embeddings, candidate sets, and simplex restriction and differ only
in the power of the displacement cost.

\begin{table}[H]
  \centering
  \caption{Factor-conditioned transport-order sensitivity. Excess returns are
    conditioned on firm-specific intercepts and loadings on the Fama--French
    five factors and momentum. The 95\% intervals use 2,000 paired
    joint-date stationary-bootstrap refits. The held-out score fits all
    parameters on 2018--2022 returns and evaluates the mean Gaussian
    quasi-log score per firm-day on 2023--2026 returns; higher is better.
  The adjustment index is $\hat\lambda=\hat\rho/(1-\hat\rho)$.}
  \label{tab:transport-order}
  \begin{tabular}{lrrrr}
    \toprule
    Interaction field & $\hat\rho$ & 95\% interval & $\hat\lambda$
    & Held-out score \\
    \midrule
    Target-anchored W1 &
    \ppnum[3]{transport-w1-rho} &
    $[\ppnum[3]{transport-w1-rho-ci-lower},
    \ppnum[3]{transport-w1-rho-ci-upper}]$ &
    \ppnum[2]{transport-w1-lambda} &
    \ppnum[6]{transport-w1-predictive-score} \\
    Target-anchored W2 &
    \ppnum[3]{transport-w2-rho} &
    $[\ppnum[3]{transport-w2-rho-ci-lower},
    \ppnum[3]{transport-w2-rho-ci-upper}]$ &
    \ppnum[2]{transport-w2-lambda} &
    \ppnum[6]{transport-w2-predictive-score} \\
    \bottomrule
  \end{tabular}
\end{table}

W1 yields $\hat\rho=\ppnum[3]{transport-w1-rho}$, and the paired W2-minus-W1
difference is $\ppnum[6]{transport-delta-rho}$, with 95\% interval
$[\ppnum[6]{transport-delta-rho-ci-lower},
\ppnum[6]{transport-delta-rho-ci-upper}]$.
The difference is precisely estimated but economically negligible, and W2 has
the slightly lower post-corpus fit, with a quasi-log-likelihood difference of
$\ppnum[2]{transport-delta-loglik}$
$[\ppnum[2]{transport-delta-loglik-ci-lower},
\ppnum[2]{transport-delta-loglik-ci-upper}]$.
Out of period, the W2-minus-W1 score difference is
$\ppnum[6]{transport-delta-predictive-score}$
$[\ppnum[6]{transport-delta-predictive-score-ci-lower},
\ppnum[6]{transport-delta-predictive-score-ci-upper}]$.
Its 90\% interval, from \ppnum[6]{transport-equivalence-ci-lower} to
\ppnum[6]{transport-equivalence-ci-upper}, lies inside the
$\pm\ppnum[6]{transport-equivalence-margin}$ margin fixed in advance, so
the two fields are \ppvalue{transport-equivalence-status}, and the equivalence
holds across the factor ladder.

W1 and W2 are nearly the same object: their off-diagonal weights correlate at
\ppnum[3]{transport-weight-correlation},
their active supports have Jaccard overlap
\ppnum[3]{transport-active-support-jaccard}, and their induced peer-return
signals correlate at \ppnum[4]{transport-signal-correlation}, with a difference
whose standard deviation is \ppnum[3]{transport-relative-signal-difference-sd}
of the W2 signal's.
Factor residualization leaves that classification unchanged. Both separation
thresholds bind, so the comparison rests on the paired contrasts above; a joint
W1 and W2 specification, added after the standalone comparison, places its
optimum on the W1-only boundary with strongly negatively correlated channel
estimates (\Cref{tab:transport-order-joint}).

A dimension ladder over the normalized Qwen3-Embedding-8B prefix locates the
equivalence in the construction
(\ppAppendixRef{sec:p5-matryoshka},
\ppAppendixRef{sec:app-transport-order}).
Shorter prefixes change which articles the two orders match, but the changed
assignments do not survive reconstruction: operator and spatial-lag
correlations remain effectively one at every dimension, and W2 never improves
the held-out score.
Target reconstruction therefore maps different article-level alignments into
nearly identical financial operators.

\subsection{Representation robustness}
\label{sec:results-representation}

The representation checks examine dependence on the embedding model and
its reported training-information vintage.
Across Qwen3-Embedding widths from 4,096 to 64 coordinates, the smaller
Qwen3-Embedding-4B model, the BGE-large-v1.5 model family, and all eight
published ChronoBERT checkpoints with reported 2017--2024 training cutoffs,
the barycentric field attains a higher quasi-log-likelihood than its matching
RBF field in the raw-return, common-intercept specification, while both fields
retain positive spatial coefficients
(\Cref{tab:p5-representation-ablation}).
Magnitudes respond to aggressive compression: from 1,024 coordinates downward,
barycentric feedback and its margin over RBF decline.
The annual ChronoBERT staircase is non-monotone: many paired changes in
barycentric feedback and its gap relative to RBF reject zero after Holm
adjustment, while only a subset meets the post-specified joint equivalence
rule (\Cref{tab:p5-representation-contrasts}).
The first annual step rejects zero for both estimands and meets joint
equivalence; the 2017-to-2024 endpoint meets joint equivalence without
Holm-adjusted rejection for either estimand.
These cumulative-training comparisons do not identify a causal leakage effect.
The pre-period ChronoBERT 2017 check preserves a positive barycentric gain of
\ppnum[4]{point-in-time-w2-held-out-gain} and a barycentric-minus-RBF
difference of
\ppnum[4]{point-in-time-w2-minus-rbf-estimate}, with 95\% interval
$[\ppnum[4]{point-in-time-w2-minus-rbf-ci-lower},
\ppnum[4]{point-in-time-w2-minus-rbf-ci-upper}]$
(\ppAppendixRef{sec:p5-point-in-time}).
Annual re-estimation under the frozen field yields positive adjustment in each
of 2023, 2024, 2025, and the partial 2026 sample, for the barycentric field
alone and jointly with the news-link field
(\Cref{tab:qmle-annual,tab:two-field-annual}).
